\subsection{实形式}

若 \(\mathfrak{a}\) 是复李代数，我们用 \(\mathfrak{a}_{[\mathbf{R} ]}\) （有时也简记为 \(\mathfrak{a}\) ）表示由标量限制得到的实李代数。若 \(\mathfrak{g}\) 是实李代数，我们用 \(\mathfrak{g}_{(\mathbf{C})}\) （有时也简记为 \(\mathfrak{g}_{\mathbf{C}}\) ）表示由标量扩张得到的复李代数 \(\mathbf{C}\otimes_{\mathbf{R}}\mathfrak{g}\) 。实李代数同态 \(\mathfrak{g}\to \mathfrak{a}_{[\mathbf{R}]}\) 与复李代数同态 \(\mathfrak{g}_{(\mathbf{C})}\to \mathfrak{a}\) 成一一对应：若 \(f:\mathfrak{g}\to \mathfrak{a}_{[\mathbf{R}]}\) 与 \(g:\mathfrak{g}_{(\mathbf{C})}\to \mathfrak{a}\) 相对应，则 \(f(x) = g(1\otimes x)\) 且 \(g(\lambda \otimes x) = \lambda f(x)\) 对 \(x\in \mathfrak{g},\lambda \in \mathbf{C}\) 成立。

定义 1. 设 \(\mathfrak{a}\) 是复李代数。\(\mathfrak{a}\) 的实形式是指实子代数 \(\mathfrak{g}\) （\(\mathfrak{a}\) 的），它是一个 \(\mathbf{R}\) -结构（在 \(\mathbf{C}\) -向量空间 \(\mathfrak{a}\) 上）（《代数学》，第 II 章，§8，第 1 节，定义 1）。

这就是说，复李代数同态 \(\mathfrak{g}_{(\mathbf{C})}\to\mathfrak{a}\) （与典则单射 \(g\to a_{[R]}\) 相伴）是双射。于是，a 的实子代数 g 是 a 的实形式，当且仅当实向量空间 a 的子空间 g 与 ig 互为余补。a 相对于实形式 g 的共轭是满足下式的映射 \(\sigma:a\to a\)

\[
\sigma (x + i y) = x - i y, \quad x, y \in \mathfrak {g}.\tag{1}
\]

命题 1. a) 设 \(\mathfrak{g}\) 是 \(\mathfrak{a}\) 的实形式，\(\sigma\) 是 \(\mathfrak{a}\) 相对于 \(\mathfrak{g}\) 的共轭。则：

\[
\sigma^ {2} = \mathrm{Id} _ {\mathfrak {a}}, \quad \sigma (\lambda x + \mu y) = \bar {\lambda} \sigma (x) + \bar {\mu} \sigma (y), \quad [ \sigma (x), \sigma (y) ] = \sigma [ x, y ]\tag{2}
\]

其中 \(\lambda, \mu \in \mathbf{C}\) ，\(x, y \in \mathfrak{a}\) 。元素 \(x\) （\(\mathfrak{a}\) 中的）属于 \(\mathfrak{g}\) 当且仅当 \(\sigma(x) = x\) 。b) 设 \(\sigma: \mathfrak{a} \to \mathfrak{a}\) 是满足 (2) 的映射。则不动点所成的集 \(\mathfrak{g}\) （\(\sigma\) 的）是 \(\mathfrak{a}\) 的实形式，且 \(\sigma\) 是 \(\mathfrak{a}\) 相对于 \(\mathfrak{g}\) 的共轭。

证明是直接的。

注意，若 B 表示 a 的 Killing 型，且 g 是 a 的实形式，则 B 在 g 上的限制是 g 的 Killing 型；特别地，B 在 \(g \times g\) 上取实值。假定 a 约化；则实李代数 g 紧当且仅当 B 在 g 上的限制是负的（§1 第 3 节）。这时我们称 g 是 a 的紧实形式。

